% !TeX program = xelatex
\documentclass[a4paper,10pt]{article}

% ============================================================================
% សន្លឹកកិច្ចការសិស្ស ស-ខ្មៅ (Student Worksheet - Grade 8 2.4.6c ROOT2)
% បរិបទ៖ គ្មានស្លាយ LCD និងគ្មានម៉ាស៊ីនគិតលេខ — សិស្សបំពេញផ្ទាល់លើសន្លឹកកិច្ចការ ២ ទំព័រ A4
% ============================================================================
\usepackage[a4paper, left=0.85cm, right=0.85cm, top=0.70cm, bottom=0.95cm, footskip=0.48cm]{geometry}
\usepackage{amsmath, amssymb}
\usepackage{unicode-math}
\usepackage{fontspec}
\usepackage{array, colortbl}
\usepackage[table]{xcolor}
\usepackage{enumitem}
\usepackage{pifont}
\usepackage{fancyhdr}
\usepackage{lastpage}
\usepackage{tikz}
\usepackage[most]{tcolorbox}
\usetikzlibrary{calc, arrows.meta, decorations.pathreplacing, positioning, babel}

\makeatletter
\let\@expl@@@mark@update@singlecol@structures@@\relax
\makeatother

\providecommand{\NGSFindFonts}{%
  \IfFileExists{C:/Windows/Fonts/cambria.ttc}{%
    \def\NGSSysPath{C:/Windows/Fonts/}%
    \def\NGSUserPath{C:/Users/KIT/AppData/Local/Microsoft/Windows/Fonts/}%
  }{%
    \IfFileExists{./fonts/KhmerOSsiemreap.ttf}{\def\NGSSysPath{./fonts/}\def\NGSUserPath{./fonts/}}{%
    \IfFileExists{../fonts/KhmerOSsiemreap.ttf}{\def\NGSSysPath{../fonts/}\def\NGSUserPath{../fonts/}}{%
    \IfFileExists{../../fonts/KhmerOSsiemreap.ttf}{\def\NGSSysPath{../../fonts/}\def\NGSUserPath{../../fonts/}}{%
    \IfFileExists{../../../fonts/KhmerOSsiemreap.ttf}{\def\NGSSysPath{../../../fonts/}\def\NGSUserPath{../../../fonts/}}{%
    \IfFileExists{../../../../fonts/KhmerOSsiemreap.ttf}{\def\NGSSysPath{../../../../fonts/}\def\NGSUserPath{../../../../fonts/}}{%
    \IfFileExists{../../../../../fonts/KhmerOSsiemreap.ttf}{\def\NGSSysPath{../../../../../fonts/}\def\NGSUserPath{../../../../../fonts/}}{%
    \IfFileExists{../../../../../../fonts/KhmerOSsiemreap.ttf}{\def\NGSSysPath{../../../../../../fonts/}\def\NGSUserPath{../../../../../../fonts/}}{%
      \def\NGSSysPath{../../../../../../../fonts/}\def\NGSUserPath{../../../../../../../fonts/}%
    }}}}}}}%
  }%
}
\NGSFindFonts

\setmathfont[Path={\NGSSysPath}, FontIndex=1, Scale=1.02]{cambria.ttc}
\defaultfontfeatures{Ligatures=TeX}

\setmainfont[
  Path={\NGSUserPath}, Extension=.ttf, Script=Khmer,
  BoldFont={KhmerOSsiemreap}, BoldFeatures={FakeBold=1.8},
  ItalicFont={KhmerOSsiemreap}, ItalicFeatures={FakeSlant=0.22}
]{KhmerOSsiemreap}

\newfontfamily\muol[Path={\NGSUserPath}, Extension=.ttf, Script=Khmer]{KhmerOSmuollight}
\newfontfamily\TimesEn[Path={\NGSSysPath}, Extension=.ttf, BoldFont={timesbd}, ItalicFont={timesi}]{times}

\XeTeXlinebreaklocale "km"
\XeTeXlinebreakskip = 0pt plus 1pt minus 0.1pt

\definecolor{panelgray}{gray}{0.93}
\definecolor{headgray}{gray}{0.85}

\pagestyle{fancy}
\fancyhf{}
\renewcommand{\headrulewidth}{0pt}
\fancyfoot[L]{\fontsize{7.5pt}{9.5pt}\selectfont អនុវិទ្យាល័យព្រះក្រសាំង \en{|} គណិតវិទ្យាថ្នាក់ទី៨ \en{|} មេរៀនទី២៖ ៤.៦.គ. ការគណនាឫសការេ និងសំយោគមេរៀនទី២ (គ្រូ៖ ហុក ប៊ុនថែន)}
\fancyfoot[R]{\fontsize{7.8pt}{9.5pt}\selectfont ទំព័រ \en{\thepage/\pageref{LastPage}}}

\setlength{\parindent}{0pt}
\setlength{\parskip}{0pt}
\setlength{\tabcolsep}{3.8pt}

\newcommand{\en}[1]{{\TimesEn #1}}
\newcommand{\enbf}[1]{{\TimesEn\bfseries #1}}
\newcommand{\unit}[1]{\,\text{\en{#1}}}
\newcommand{\cnum}[1]{\mbox{\fontsize{10pt}{12pt}\selectfont\ding{\numexpr181+#1\relax}}}
\newcommand{\chkbox}{\tikz[baseline=-0.12ex]\draw[black, line width=0.8pt] (0,0) rectangle (0.24cm,0.24cm);}
\newcommand{\blank}[1][1.5cm]{\makebox[#1]{\dotfill}}
\newcommand{\ansbox}[1][0.85cm]{\,\tikz[baseline=-0.5ex]\draw[black, line width=0.65pt, fill=white] (0,0) rectangle (#1, 0.42cm);\,}
\newcommand{\secbar}[2]{%
  \begin{tcolorbox}[enhanced, sharp corners, colback=headgray, colframe=black, boxrule=0.9pt, boxsep=1pt, top=2pt, bottom=2pt, left=5pt, right=5pt, before skip=4pt, after skip=3pt]
    {\muol\fontsize{9.2pt}{13pt}\selectfont #1}\hfill{\fontsize{8pt}{11pt}\selectfont\textbf{#2}}
  \end{tcolorbox}%
}

\newcolumntype{L}[1]{>{\raggedright\arraybackslash}p{#1}}
\newcolumntype{C}[1]{>{\centering\arraybackslash}p{#1}}

\begin{document}
\fontsize{9.6pt}{14.4pt}\selectfont

% ============================================================================
% ទំព័រទី ១៖ E1 + E2 (POE 1 & POE 2) + E3 សំយោគក្បួន និងធនាគារហេតុផល
% ============================================================================
\begin{tcolorbox}[enhanced, sharp corners, colback=white, colframe=black, boxrule=1.0pt, boxsep=2pt, top=3pt, bottom=3pt, left=6pt, right=6pt]
  \begin{center}
    {\muol\fontsize{10.2pt}{15pt}\selectfont អនុវិទ្យាល័យព្រះក្រសាំង \enbf{--} សន្លឹកកិច្ចការសិស្ស \enbf{(}មេរៀនទី២៖ ៤\enbf{.}៦\enbf{.}គ\enbf{.} ការគណនាឫសការេដោយដៃ\enbf{)}}\\[3pt]
    {\fontsize{8.8pt}{12pt}\selectfont
    ឈ្មោះសិស្ស / ដៃគូ៖ \blank[5.2cm] \hspace{0.25cm}
    ថ្នាក់ទី៨៖ \blank[1.1cm] \hspace{0.25cm}
    ក្រុមទី៖ \blank[1.0cm] \hspace{0.25cm}
    កាលបរិច្ឆេទ៖ ២០/១១/២០២៦}
  \end{center}
\end{tcolorbox}

% --- E1 + E2 ---
\secbar{ផ្នែកទី១៖ បេសកកម្មទស្សន៍ទាយ \enbf{(E1)} និងការរុករកជាដៃគូ \enbf{(E2 -- POE 1 \& POE 2)}}{រយៈពេល ១៤ នាទី}

\begin{tcolorbox}[enhanced, sharp corners, colback=panelgray, colframe=black, boxrule=0.75pt, boxsep=1.5pt, top=2.5pt, bottom=2.5pt, left=5pt, right=5pt]
  \textbf{បេសកកម្ម \en{E1} (គូសធីក \chkbox\ មុនពិសោធន៍ — គ្មានម៉ាស៊ីនគិតលេខ!)៖}
  (១) ដោយ $12^2=144$ និង $13^2=169$ តើ $\sqrt{153.76}$ ស្មើ៖ \chkbox\ $124$ \quad \chkbox\ $12.4$\\
  (២) រវាង $2304$ និង $2312$ តើមួយណាជាការេប្រាកដ ($\sqrt{\phantom{x}}=48$)៖ \chkbox\ $2304$ \quad \chkbox\ $2312$
\end{tcolorbox}

\vspace{3pt}
\textbf{ក. សកម្មភាពទី ១ (\en{POE 1})៖ របៀបទី ១ — គណនាឫសការេតាមការបំបែកកត្តាបឋមជាគូ (ទំ.២៤)}

{\renewcommand{\arraystretch}{1.48}
\begin{tabular}{|C{2.45cm}|L{6.45cm}|L{5.35cm}|C{3.35cm}|}
  \hline
  \rowcolor{headgray}
  \textbf{កន្សោម} & \textbf{បំបែកជាផលគុណកត្តាបឋម} & \textbf{ផ្គុំជាគូការេ $\sqrt{a^2\times b^2} = a\times b$} & \textbf{ឫសការេ} \\
  \hline
  $\sqrt{225}$ & $225 = 9\times 25 = 3^2 \times 5^2$ & $= \sqrt{(3\times 5)^2} = 3\times 5$ & $= 15$ \\
  \hline
  $\sqrt{625}$ & $625 = 25\times 25 = 5^2 \times \blank[1.1cm]$ & $= 5 \times \blank[1.1cm]$ & $= \blank[1.5cm]$ \\
  \hline
  $\sqrt{441}$ & $441 = 9\times 49 = (\blank[0.8cm])^2 \times (\blank[0.8cm])^2$ & $= \blank[0.9cm] \times \blank[0.9cm]$ & $= \blank[1.5cm]$ \\
  \hline
  $\sqrt{784}$ & $784 = 16\times 49 = 2^4 \times 7^2 = (2^2)^2 \times 7^2$ & $= \blank[0.9cm] \times \blank[0.9cm]$ & $= \blank[1.5cm]$ \\
  \hline
  $\sqrt{900}$ & $900 = 9\times 100 = 3^2 \times 10^2$ & $= \blank[0.9cm] \times \blank[0.9cm]$ & $= \blank[1.5cm]$ \\
  \hline
  $\sqrt{1225}$ & $1225 = 25\times 49 = 5^2 \times (\blank[0.8cm])^2$ & $= \blank[0.9cm] \times \blank[0.9cm]$ & $= \blank[1.5cm]$ \\
  \hline
\end{tabular}}

\vspace{3pt}
$\bullet$ \textbf{ពន្យល់ (\en{E})៖} បើបំបែក $2312 = 8\times 289 = 2^3 \times 17^2$ មានកត្តា $2^3$ (និទស្សន្តសេស) តើ $2312$ ជាការេប្រាកដដែរឬទេ? \blank[2.1cm]

\vspace{4pt}
\textbf{ខ. សកម្មភាពទី ២ (\en{POE 2})៖ របៀបទី ២ — វិធីចែកផ្ដាច់ពីរខ្ទង់ចេញពីស្ដាំទៅឆ្វេង (អនុវត្ត ១ ទំ.២៤)}

\begin{tabular}{|L{9.20cm}|L{9.20cm}|}
  \hline
  \textbf{គំរូទី ១៖ គណនា $\sqrt{225}$ (ផ្ដាច់ពីរខ្ទង់៖ $2\,|\,25$)}\par\vspace{2pt}
  $\bullet$ \textbf{ជំហាន ១៖} រកការេធំបំផុត $\le 2$ គឺ $1^2 = 1$ ។ សរសេរឫសខ្ទង់ទី១ គឺ $\mathbf{1}$ រួចដក $2 - 1 = 1$ ។\par\vspace{2pt}
  $\bullet$ \textbf{ជំហាន ២៖} ទម្លាក់ក្រុមពីរខ្ទង់ $25$ ចុះមក បាន $\mathbf{125}$ ។\par\vspace{2pt}
  $\bullet$ \textbf{ជំហាន ៣៖} យកឫស $1\times 2 = 2$ រួចដាក់លេខ $b=\blank[0.7cm]$ ពីក្រោយ៖ $\overline{25}\times \blank[0.7cm] = 125$ (សល់ $0$) $\implies \sqrt{225} = \mathbf{\blank[1.0cm]}$ ។
  &
  \textbf{អនុវត្តផ្ទាល់ដៃ៖ គណនា $\sqrt{625}$ (ផ្ដាច់ពីរខ្ទង់៖ $6\,|\,25$)}\par\vspace{2pt}
  $\bullet$ \textbf{ជំហាន ១៖} ការេធំបំផុត $\le 6$ គឺ $(\blank[0.7cm])^2 = 4$ ។ សរសេរឫសខ្ទង់ទី១ គឺ $\blank[0.7cm]$ រួចដក $6 - 4 = \blank[0.7cm]$ ។\par\vspace{2pt}
  $\bullet$ \textbf{ជំហាន ២៖} ទម្លាក់ក្រុមពីរខ្ទង់ $25$ ចុះមក បាន $\blank[1.1cm]$ ។\par\vspace{2pt}
  $\bullet$ \textbf{ជំហាន ៣៖} យកឫស $2\times 2 = 4$ រួចដាក់លេខ $b=\blank[0.7cm]$ ពីក្រោយ៖ $\overline{4\blank[0.5cm]}\times \blank[0.6cm] = 225 \implies \sqrt{625} = \mathbf{\blank[1.0cm]}$ ។ \\
  \hline
\end{tabular}

\vspace{3pt}
$\bullet$ \textbf{ពន្យល់ (\en{E})៖} ហេតុអ្វីក្នុងជំហានទី ៣ យើងត្រូវយកឫសរកឃើញគុណនឹង $2$ រួចដាក់លេខ $b$ ពីក្រោយ? (ភ្ជាប់រូបមន្ត $(10a+b)^2 = 100a^2 + (20a+b)b$៖ មេគុណ $20a$ បានមកពី $\blank[1.4cm]\times a$) ។

% --- E3 ---
\secbar{ផ្នែកទី២៖ សំយោគក្បួនគណនា \enbf{(E3)} និងធនាគារហេតុផល \enbf{B\&W} \cnum{1}\enbf{--}\cnum{5}}{សៀវភៅពុម្ព ទំ.២៤–២៥}

\begin{tabular}{|L{9.20cm}|L{9.20cm}|}
  \hline
  \textbf{១. ក្បួនគណនាឫសការេដោយដៃ (ទំ.២៤)៖}\par\vspace{2pt}
  $\bullet$ \textbf{របៀបទី ១៖} បំបែករ៉ាឌីកង់ជាផលគុណកត្តាបឋមមាននិទស្សន្តគូ៖ $\sqrt{a^2\times b^2} = \blank[0.8cm]\times \blank[0.8cm]$ ។\par\vspace{2pt}
  $\bullet$ \textbf{របៀបទី ២៖} ផ្ដាច់ពីរខ្ទង់ម្ដងចេញពី\textbf{ក្បៀសទសភាគ} (ទៅឆ្វេង និងទៅស្ដាំ) ហើយគ្រប់ជំហានត្រូវយកឫសរកឃើញគុណនឹង $\ansbox[0.7cm]$ ។\par\vspace{2pt}
  $\bullet$ \textbf{សំយោគស្វ័យគុណ (ទំ.២៥)៖} $a^m\cdot a^n = a^{\blank[0.8cm]}$ ; $\frac{a^m}{a^n} = a^{\blank[0.8cm]}$ ; $(a^m)^n = a^{\blank[0.8cm]}$ ; $A\times 10^n$ ($1\le A < \blank[0.7cm]$) ។
  &
  \textbf{២. ធនាគារហេតុផល \en{B\&W} (\cnum{1}--\cnum{5})៖}\par\vspace{2pt}
  \cnum{1} កត្តាបឋមជាគូ៖ $\sqrt{p^{2k} q^{2m}} = p^k q^m$ (សល់និទស្សន្តសេស $\to$ មិនដាច់) ។\par\vspace{2pt}
  \cnum{2} ផ្ដាច់ពីរខ្ទង់៖ ផ្ដាច់ពីរខ្ទង់ម្ដងចេញពីក្បៀសទសភាគ ។\par\vspace{2pt}
  \cnum{3} គុណឫសនឹង $2$៖ រកលេខ $b$ ដែល $\overline{(2a)b}\times b \le$ សំណល់ទម្លាក់ ។\par\vspace{2pt}
  \cnum{4} របងប៉ាន់ស្មាន៖ $k^2 < N < (k+1)^2 \implies k < \sqrt{N} < k+1$ ។\par\vspace{2pt}
  \cnum{5} ប្រភាគ/ទសភាគ៖ $\sqrt{\frac{a}{b}} = \frac{\sqrt{a}}{\sqrt{b}}$ ; ខ្ទង់ក្រោយក្បៀសនៅសល់ពាក់កណ្ដាល ។ \\
  \hline
\end{tabular}

\vspace{4pt}
{\renewcommand{\arraystretch}{1.48}
\begin{tabular}{|C{4.40cm}|C{4.40cm}|C{4.40cm}|C{4.40cm}|}
  \hline
  \rowcolor{panelgray}
  $\sqrt{1.44} = \blank[1.7cm]$ & $\sqrt{0.0169} = \blank[1.5cm]$ & $\sqrt{\dfrac{81}{100}} = \blank[1.6cm]$ & $\sqrt{3^2\times 7^2} = \blank[1.4cm]$ \\
  \hline
  $\sqrt{2.25} = \blank[1.7cm]$ & $\sqrt{0.0036} = \blank[1.5cm]$ & $\sqrt{\dfrac{121}{144}} = \blank[1.5cm]$ & $\sqrt{2^4\times 5^2} = \blank[1.4cm]$ \\
  \hline
  $\sqrt{4.84} = \blank[1.7cm]$ & $\sqrt{0.0625} = \blank[1.5cm]$ & $\sqrt{\dfrac{169}{225}} = \blank[1.5cm]$ & $\sqrt{3^4\times 2^2} = \blank[1.4cm]$ \\
  \hline
\end{tabular}}

\newpage

% ============================================================================
% ទំព័រទី ២៖ E4 លំហាត់គំរូ + E5 ប្រតិបត្តិ + វិនិច្ឆ័យ A-D + Exit Check + កិច្ចការផ្ទះ
% ============================================================================
\secbar{ផ្នែកទី៣៖ ពង្រីកចំណេះដឹងតាមជំហានស៊ើបអង្កេត \enbf{(E4 --} អនុវត្ត ២ ទំ\enbf{.}២៤\enbf{)}}{រយៈពេល ១០ នាទី}

\textbf{១. គណនាឫសការេដោយដៃតាមវិធីចែកផ្ដាច់ពីរខ្ទង់ និងផ្ទៀងដោយរបងប៉ាន់ស្មានទំហំ (អនុវត្ត ២ និងប្រតិបត្តិ ១ ទំ.២៤)៖}

{\renewcommand{\arraystretch}{1.48}
\begin{tabular}{|L{3.35cm}|L{9.15cm}|C{3.10cm}|C{2.20cm}|}
  \hline
  \rowcolor{headgray}
  \textbf{លំហាត់} & \textbf{ជំហានផ្ដាច់ពីរខ្ទង់ និងគុណឫសនឹង $2$} & \textbf{ចម្លើយ} & \textbf{ហេតុផល} \\
  \hline
  \textbf{ក.} $\sqrt{529}$ & ផ្ដាច់ $5\,|\,29$៖ $2^2=4$ (សល់ $129$) ; $2\!\times\! 2=4 \to \overline{43}\times \blank[0.6cm] = 129$ & $= \blank[2.0cm]$ & \cnum{2}, \cnum{3} \\
  \hline
  \textbf{ខ.} $\sqrt{153.76}$ & ផ្ដាច់ $1\,|\,53\,.\,76$៖ បាន $12$ សល់ $9$ ទម្លាក់ $76$ បាន $976$ ; $\overline{24\blank[0.5cm]}\times 4 = 976$ & $= \blank[2.0cm]$ & \cnum{2}, \cnum{4} \\
  \hline
  \textbf{គ.} $\sqrt{2304}$ & ផ្ដាច់ $23\,|\,04$៖ $4^2=16$ (សល់ $704$) ; $4\!\times\! 2=8 \to \overline{88}\times \blank[0.6cm] = 704$ & $= \blank[2.0cm]$ & \cnum{2}, \cnum{3} \\
  \hline
\end{tabular}}

\vspace{4pt}
\begin{tcolorbox}[enhanced, sharp corners, colback=panelgray, colframe=black, boxrule=0.75pt, boxsep=1.5pt, top=2.5pt, bottom=2.5pt, left=5pt, right=5pt]
  \textbf{កំណត់ចំណាំកែតម្រូវពុម្ព ទំ.២៤៖}
  (១) ចំពោះ $\sqrt{153.76}$៖ ដោយ $12^2=144 < 153.76 < 169=13^2$ ចម្លើយត្រឹមត្រូវគឺ $\mathbf{12.4}$ (មិនមែន $124$ ទេ!) ។\\
  (២) ចំពោះ $2312$៖ $2312 - 48^2 = 2312 - 2304 = \blank[0.9cm]$ (សល់សំណល់ $\ne 0$ ព្រោះ $2312 = 2^3\times 17^2$ មិនមែនការេប្រាកដទេ!) ។
\end{tcolorbox}

\vspace{3pt}
$\bullet$ \textbf{សំណួរឆ្លុះបញ្ចាំង \en{(EIA Looking-Back)}៖} ទសភាគ $0.208849$ មាន $6$ ខ្ទង់ក្រោយក្បៀស តើ $\sqrt{0.208849}$ មានប៉ុន្មានខ្ទង់ក្រោយក្បៀស? \blank[1.8cm]

% --- E5 ---
\secbar{ផ្នែកទី៤៖ វាយតម្លៃលទ្ធផលសិក្សា \enbf{(E5 --} ប្រតិបត្តិ ២ ទំ\enbf{.}២៤ \enbf{+} វិនិច្ឆ័យ \enbf{A--D +} សំបុត្រចេញ\enbf{)}}{រយៈពេល ១០ នាទី}

\textbf{ក. លំហាត់ប្រតិបត្តិ ២ (សៀវភៅពុម្ព ទំ.២៤ — គណនាដោយដៃជាដៃគូ រួចប្ដូរសន្លឹកផ្ទៀង \en{Pause-Review})៖}

{\renewcommand{\arraystretch}{1.52}
\begin{tabular}{|L{8.90cm}|L{8.90cm}|}
  \hline
  \textbf{ក.} $\sqrt{3249}$ ($32\,|\,49$)៖ $5^2=25$ (សល់ $749$), $\overline{10\blank[0.5cm]}\times 7=749 \Rightarrow \blank[1.4cm]$ &
  \textbf{ខ.} $\sqrt{\dfrac{2025}{169}} = \dfrac{\sqrt{2025}}{\sqrt{169}} = \dfrac{\blank[1.1cm]}{\blank[1.1cm]}$ \quad (\cnum{5}) \\
  \hline
  \multicolumn{2}{|L{18.0cm}|}{\textbf{គ.} $\sqrt{0.208849}$ ($0.\,20\,|\,88\,|\,49$)៖ $4^2=16$ (សល់ $488$), $\overline{85}\times 5=425$ (សល់ $6349$), $\overline{907}\times 7=6349 \implies \sqrt{0.208849} = \blank[2.0cm]$ (\cnum{5})} \\
  \hline
  \textbf{ឃ. (ពង្រឹង \#16)៖} $\sqrt{1225} = \sqrt{5^2\times 7^2} = \blank[2.2cm]$ (\cnum{1}) &
  \textbf{ង. (ពង្រឹង \#17)៖} $\sqrt{0.0625} = \blank[2.4cm]$ (\cnum{5}) \\
  \hline
\end{tabular}}

\vspace{5pt}
\textbf{ខ. គូសវង់ជុំវិញជម្រើសត្រឹមត្រូវតែ ១ គត់ ក្នុងចំណោម \en{A, B, C, D} រួចសរសេរលេខហេតុផល \cnum{1}--\cnum{5} នៃជម្រើសខុស៖}

{\renewcommand{\arraystretch}{1.35}
\begin{tabular}{|C{4.40cm}|C{4.40cm}|C{4.40cm}|C{4.40cm}|}
  \hline
  \cellcolor{panelgray}\textbf{ជម្រើស \en{A}} & \cellcolor{panelgray}\textbf{ជម្រើស \en{B}} & \cellcolor{panelgray}\textbf{ជម្រើស \en{C}} & \cellcolor{panelgray}\textbf{ជម្រើស \en{D}} \\
  \hline
  $\sqrt{0.208849} = 0.457$ & $\sqrt{153.76} = 124$ & $\sqrt{2312} = 48$ ដាច់គត់ & $\sqrt{0.09} = 0.03$ \\[3pt]
  \chkbox\ ត្រូវ \quad \chkbox\ ខុស (\blank[0.7cm]) & \chkbox\ ត្រូវ \quad \chkbox\ ខុស (\blank[0.7cm]) & \chkbox\ ត្រូវ \quad \chkbox\ ខុស (\blank[0.7cm]) & \chkbox\ ត្រូវ \quad \chkbox\ ខុស (\blank[0.7cm]) \\
  \hline
\end{tabular}}

\vspace{5pt}
\begin{tcolorbox}[enhanced, sharp corners, colback=white, colframe=black, boxrule=0.85pt, boxsep=2.5pt, top=3.5pt, bottom=3.5pt, left=5pt, right=5pt]
  {\muol\fontsize{8.8pt}{12.5pt}\selectfont សំបុត្រចេញសំយោគមេរៀនទី២ \enbf{(Exit Check --} ៦០ វិនាទី\enbf{)}៖}\\[2pt]
  (១) $\sqrt{625} = \blank[1.6cm]$ \quad \en{|} \quad
  (២) $\sqrt{1.44} = \blank[1.6cm]$ \quad \en{|} \quad
  (៣) ទម្រង់ស្តង់ដានៃ $0.0004 = \blank[2.2cm]$
\end{tcolorbox}

\vspace{5pt}
\begin{tcolorbox}[enhanced, sharp corners, colback=panelgray, colframe=black, boxrule=0.85pt, boxsep=2.5pt, top=3.5pt, bottom=3.5pt, left=5pt, right=5pt]
  {\muol\fontsize{8.8pt}{12.5pt}\selectfont កិច្ចការផ្ទះ \enbf{(}សៀវភៅពុម្ព ទំព័រ ២៦ លំហាត់ \enbf{\#16} និង \enbf{\#17 +} រំលឹកសង្ខេបទំ\enbf{.}២៥\enbf{)}៖}\\[2pt]
  $\bullet$ \textbf{លំហាត់ \#16៖} គណនាឫសការេដោយដៃ៖ ក) $\sqrt{441}$ ; ខ) $\sqrt{784}$ ; គ) $\sqrt{1225}$ ; ឃ) $\sqrt{1764}$ ។\\
  $\bullet$ \textbf{លំហាត់ \#17៖} គណនាឫសការេនៃទសភាគ និងប្រភាគក្នុងសៀវភៅពុម្ព ទំ.២៦ ព្រមទាំងអានមុន \textbf{មេរៀនទី៣៖ ទំហំសមាមាត្រស្រប (ទំ.២៧)} ។
\end{tcolorbox}

\end{document}
